% DOCUMENT COOKBOOK / SHARED SETUP. INPUT THIS FILE; DO NOT COMPILE IT.
% The ordinary document remains flowing text. Every optional feature below
% stays inactive unless explicitly used in a document or recipe.

% FOUNDATION: page geometry and typography
\usepackage[T1]{fontenc}
\usepackage{lmodern,microtype}
\usepackage[letterpaper,margin=1in]{geometry} % change paper/margins here
\usepackage{setspace}
\setlength{\parindent}{1.3em}
\setlength{\parskip}{0pt}

% Scientific writing and layout
\usepackage{amsmath,amssymb,amsthm,bm}
\usepackage{graphicx,caption,subcaption,wrapfig,float,placeins}
\usepackage{booktabs,tabularx,array,multicol}
\usepackage{xcolor,tikz}
\usetikzlibrary{calc,positioning}
\usepackage{titlesec}
\usepackage{fancyhdr}
\usepackage{listings}
\definecolor{CUGold}{HTML}{CFB87C}
\definecolor{Ink}{HTML}{193848}
\definecolor{Teal}{HTML}{237C89}
\definecolor{Muted}{HTML}{5E7179}
\definecolor{Soft}{HTML}{EFF5F6}
\graphicspath{{figures/}}
\captionsetup{font=small,labelfont=bf}
\newtheorem{theorem}{Theorem}[section]
\theoremstyle{definition}
\newtheorem{definition}[theorem]{Definition}
\lstset{basicstyle=\ttfamily\footnotesize,frame=single,breaklines=true,
  columns=fullflexible,keepspaces=true,showstringspaces=false,
  keywordstyle=\color{Teal},commentstyle=\color{Muted}}

% Citation system: biblatex + Biber. Do not mix with natbib/BibTeX commands.
\usepackage[backend=biber,style=numeric,sorting=none]{biblatex}
\addbibresource{references.bib}
\usepackage[colorlinks=true,linkcolor=Ink,citecolor=Teal,urlcolor=Teal]{hyperref}
\usepackage[noabbrev,nameinlink]{cleveref}

% FOUNDATION: customize an OPTIONAL precision-placed title page.
% First fontsize argument is the font size (pt); second is baseline spacing.
\newcommand{\DocTitleFont}{\color{Ink}\fontsize{25}{31}\selectfont\bfseries}
\newcommand{\DocSubtitleFont}{\color{Muted}\fontsize{14}{20}\selectfont}
\newcommand{\DocAuthorFont}{\color{Ink}\fontsize{13}{18}\selectfont}
\newcommand{\DocAffiliationFont}{\color{Muted}\fontsize{11}{15}\selectfont}

% PAGE coordinates: X from page left, Y downward from page top (cm).
% Page objects are overlays: they occupy NO flow space; use on title pages.
% \DocPageTextAt{X}{Y}{WIDTH}{TEXT}
\newcommand{\DocPageTextAt}[4]{%
  \begin{tikzpicture}[remember picture,overlay]
    \node[anchor=north west,inner sep=0pt,outer sep=0pt]
      at ([xshift=#1cm,yshift=-#2cm]current page.north west) {%
        \begin{minipage}[t]{#3}\vspace{0pt}#4\par\end{minipage}};
  \end{tikzpicture}%
}
% \DocPageImageAt{X}{Y}{WIDTH}{FILE}
\newcommand{\DocPageImageAt}[4]{%
  \begin{tikzpicture}[remember picture,overlay]
    \node[anchor=north west,inner sep=0pt,outer sep=0pt]
      at ([xshift=#1cm,yshift=-#2cm]current page.north west)
      {\includegraphics[width=#3]{#4}};
  \end{tikzpicture}%
}

% CONTAINER-LOCAL coordinates: X right, Y down from the local TikZ origin.
% Place these ONLY within a tikzpicture (usually inside a figure/minipage).
% \DocImageAt{X}{Y}{WIDTH}{FILE}
\newcommand{\DocImageAt}[4]{%
  \node[anchor=north west,inner sep=0pt,outer sep=0pt] at (#1,-#2)
    {\includegraphics[width=#3]{#4}};%
}
% \DocTextAt{X}{Y}{WIDTH}{TEXT}
\newcommand{\DocTextAt}[4]{%
  \node[anchor=north west,inner sep=0pt,outer sep=0pt] at (#1,-#2) {%
    \begin{minipage}[t]{#3}\vspace{0pt}#4\par\end{minipage}};%
}
% \DocCanvas{WIDTH}{HEIGHT}{CONTENT}: explicit space, cm; no clipping.
\newcommand{\DocCanvas}[3]{%
  \begin{tikzpicture}[x=1cm,y=1cm]
    \useasboundingbox (0,0) rectangle (#1,-#2);
    #3
  \end{tikzpicture}%
}
% \DocAutoCanvas{CONTENT}: auto-sized bounding box; no reserved height.
\newcommand{\DocAutoCanvas}[1]{%
  \begin{tikzpicture}[x=1cm,y=1cm]
    \path (0,0);#1
  \end{tikzpicture}%
}

% OPT-IN REPEATED FOOTER. The blank document does NOT enable this.
% \DocSetFooter{LEFT}{CENTER}{RIGHT}{RULE THICKNESS}
\newcommand{\DocSetFooter}[4]{%
  \fancyhf{}\fancyfoot[L]{\footnotesize #1}
  \fancyfoot[C]{\footnotesize #2}\fancyfoot[R]{\footnotesize #3}
  \renewcommand{\headrulewidth}{0pt}%
  \renewcommand{\footrulewidth}{#4}%
  \pagestyle{fancy}%
}
